\begingroup
\shorthandoff{"<>"}
\setlength{\arrayrulewidth}{1pt}
\arrayrulecolor{vlepsbased}
\begin{longtable}{|p{15cm}|}
  \hline
  \rowcolor{maincolor} 
  \multicolumn{1}{|l|}{\textcolor{white}{\textbf{Caso de Uso}}} \\
  \hline
  \endfirsthead
  \multicolumn{1}{c}{{\tablename\ \thetable{} -- Continuación}} \\
  \hline
  \rowcolor{maincolor} 
  \multicolumn{1}{|l|}{\textcolor{white}{\textbf{Caso de Uso (Continuación)}}} \\
  \hline
  \endhead
  \hline \multicolumn{1}{|r|}{\textit{Continúa en la siguiente página...}} \\ \hline
  \endfoot
  \hline
  \caption{Caso de Uso CU12 — Consulta del historial de sesiones y mensajes.}
  \label{TB:CU12_HISTORIAL}
  \endlastfoot

  \rowcolor{ulepsbased}
  \begin{tabularx}{\linewidth}{@{}X X@{}}
    \textbf{Identificador:} CU12 & \textbf{Actor principal:} Usuario
  \end{tabularx} \\ \hline

  \textbf{Nombre:} Consulta del historial de sesiones y mensajes \\ \hline

  \rowcolor{ulepsbased}
  \textbf{Descripción:} \newline 
  El usuario accede a la plataforma y recupera automáticamente la lista de sus sesiones pasadas. Al seleccionar una, el sistema recupera todos los mensajes asociados y los visualiza cronológicamente en el chat. \\ \hline

  \textbf{Precondiciones:} \newline
  1. El usuario ha iniciado sesión y existen sesiones en su cuenta. \\ \hline

  \rowcolor{ulepsbased}
  \textbf{Flujo principal:} \newline
  1. El usuario accede a la interfaz principal de la aplicación. \newline
  2. El frontend solicita el listado paginado de sesiones asociadas a la cuenta. \newline
  3. El backend valida la autenticación y retorna los registros ordenados por fecha de última actualización. \newline
  4. La interfaz pinta este historial de conversaciones en el panel lateral. \newline
  5. El usuario hace clic sobre una de las sesiones anteriores. \newline
  6. El frontend envía una petición al backend para obtener los mensajes de esa sesión concreta. \newline
  7. El backend devuelve la lista de mensajes ordenados cronológicamente (con su contenido, remitente y fecha). \newline
  8. La interfaz principal limpia la vista actual y renderiza el historial recuperado, permitiendo al usuario retomar la conversación o revisarla. \\ \hline

  \rowcolor{ulepsbased}
  \textbf{Postcondiciones:} \newline
  1. El estado activo de la interfaz se sincroniza con la sesión seleccionada, preparándose para recibir nuevos mensajes (enlaza con el CU1). \\ \hline

  \textbf{Flujos alternativos:} \newline
  Ninguno. \\ \hline

\end{longtable}
\endgroup
